\maketitle

\begin{abstract}
We study neural counterfactual-regret learning in a small trading game with hidden state. The full Figgie game is multi-player and general-sum, so we do not claim an equilibrium calculation for it. Instead, we define Figgie-Lite, a two-player zero-sum reduction with four hidden goal-suit worlds, private hands, public buy outcomes, and an exact evaluator. Each information set contains an exact Bayesian card-consistency posterior computed from the private hand, together with public action history. We implement an external-sampling Deep CFR learner with neural advantage models, reservoir memories, and a learned average strategy. In a five-seed, two-tick experiment at 2,000 iterations, its mean exact exploitability is $5.73 \pm 0.23$ payoff units, compared with $1.09 \pm 0.03$ for tabular CFR. The result is negative but informative: on this tiny game, neural function approximation does not yet beat the table it is intended to replace.
\end{abstract}

\section{Introduction}

Imperfect-information learning combines two distinct problems: infer what is hidden and choose an action that is robust to uncertainty. Figgie is a market-themed card game from Jane Street with private hands and a hidden suit assignment \citep{figgierules}. Its full form is multi-player and general-sum, which makes a standard two-player exploitability calculation inappropriate.

This report asks a narrower question. In a reduced Figgie game with an exact evaluator, what happens when a Deep CFR learner receives an explicit Bayesian card-consistency feature? The goal is not to make a broad claim about reasoning, nor to present MCTS as the main result. Search motivated the original project, but the present report treats it only as background. The main object is a neural approximation to counterfactual regret learning and a transparent comparison to tabular CFR.

\section{Related work}

Counterfactual regret minimization (CFR) minimizes counterfactual regret in two-player imperfect-information games \citep{zinkevich2007regret}. Monte Carlo CFR makes those updates practical through sampled traversals, including external sampling \citep{lanctot2009mccfr}. Deep CFR replaces tabular advantage and strategy representations with neural approximators and replay memories \citep{brown2019deepcfr}. Those results motivate our implementation, but they do not imply that function approximation helps in every small game.

\section{Figgie-Lite and Bayesian inference}

Figgie-Lite retains the hidden-world structure of Figgie while removing the parts that prevent a standard exploitability calculation. There are two players and four possible goal suits. Conditional on goal suit $g$, the deck contains three cards of suit $g$ and one card of every other suit. Each player receives three cards. At each of two ticks, both players choose \textsc{Pass} or \textsc{Buy}$(s)$ for one of four suits. A successful buy transfers one card at a fixed price and the success indicator is public. Settlement is the difference between players' cash plus ten times their goal-suit card count.

The reduction is two-player and zero-sum, so strategy value and best-response exploitability can be evaluated by enumerating all four worlds and compatible deals. It is not a model of the full market mechanics, and no result below transfers as an equilibrium guarantee to full Figgie.

Let $h$ be a player's initial private hand and $g \in \{0,1,2,3\}$ the hidden goal suit. With a uniform prior over $g$ and a uniform deal within each world, the exact root posterior is
\[
P(g\mid h) = \frac{\#\{(h,h') \in \mathrm{Splits}(g)\}}{\sum_{g'} \#\{(h,h') \in \mathrm{Splits}(g')\}}.
\]
In words, a world receives probability in proportion to the number of deals that could have produced the observed hand. A hand containing three cards of one suit identifies that suit as the goal world in Figgie-Lite. The encoder also receives the full public history: both actions and both buy-success flags at each earlier tick. This feature is exact card-consistency inference, not a full posterior over strategic behavior. The neural network must learn any additional action likelihood from history.

\section{Deep CFR implementation}

The implementation follows the advantage-and-strategy decomposition of Deep CFR \citep{brown2019deepcfr}. Each player has a small two-hidden-layer ReLU advantage network. Its input is the four private card counts, four Bayesian posterior probabilities, a player indicator, and a fixed-width encoding of public history. Positive predicted advantages are normalized by regret matching to form the current policy.

At every iteration, we sample one world and deal, then run one external-sampling traversal for each player. The traversing player enumerates its five actions while the other player's action is sampled from its current regret-matched policy. Instantaneous action-value differences are added to that player's bounded reservoir memory. Visited behavior strategies, with player reach weights, enter a separate strategy reservoir. At fixed intervals the advantage networks are fit by mean squared error; after training, the average strategy network is fit by weighted cross-entropy. All reported policies are evaluated by the exact tree evaluator, not sampled self-play.

\section{Experiment}

We compare tabular chance-sampled CFR with the neural Deep CFR approximation at the same two-tick horizon and 2,000 training iterations. The fixed random seeds are $\{3,7,11,19,23\}$. Deep CFR uses two 64-unit hidden layers, a 20,000-item memory cap, updates every 25 iterations, and Adam at $2\times10^{-3}$. The table reports total exploitability $\varepsilon_0+\varepsilon_1$ in payoff units. Lower is better. Means and standard deviations are over seeds; exploitability itself is exact for the reduced game.

\begin{table}[t]
\centering
\scriptsize
\caption{Exact Figgie-Lite exploitability after 2,000 iterations.}
\label{tab:main-result}
\begin{tabular}{lrrrrrr}
\toprule
Method & Seed 3 & Seed 7 & Seed 11 & Seed 19 & Seed 23 & Mean $\pm$ s.d. \\
\midrule
Tabular CFR & 1.122 & 1.087 & 1.110 & 1.045 & 1.085 & 1.090 $\pm$ 0.030 \\
Deep CFR & 5.698 & 5.638 & 5.710 & 6.116 & 5.487 & 5.730 $\pm$ 0.233 \\
\bottomrule
\end{tabular}
\end{table}

\section{Results and limitations}

At this budget, the Deep CFR policy is consistently more exploitable than tabular CFR (Table~\ref{tab:main-result}). This is not evidence that Deep CFR is invalid. It is evidence that neural approximation carries a cost when the information-set table is small enough to represent directly. Tabular CFR has no approximation error and reaches substantially lower exploitability with the same number of sampled chance outcomes.

The Bayesian feature separates what the game reveals directly from what a network must infer. The private hand determines an exact root consistency posterior, while public buy outcomes and actions remain available for learning behavioral signal. This experiment does not isolate a posterior-versus-no-posterior ablation, so it does not establish that the feature alone causes the measured performance. That ablation, and larger games where the tabular baseline no longer fits, are the appropriate next experiments.

MCTS remains historical inspiration for the broader Figgie project, not the central method or source of results in this manuscript. The conclusions are limited to this reduction, architecture, and training budget.

\section{Reproducibility}

The reduction is implemented in \texttt{figgie\_lite.py}; the neural learner is \texttt{deep\_cfr.py}; and the exact evaluator is \texttt{cfr.py}. The command-line runner \texttt{run\_deep\_cfr.py} recreates the fixed-seed comparison. Unit tests check the fixed-width encoder, exact card-consistency posterior, valid neural strategy probabilities, and finite exact exploitability. The fixed seeds and hyperparameters above are enough to recreate the reported table without access to the full Figgie simulator.

\bibliographystyle{plainnat}
\bibliography{references}
